\documentclass[10pt,a4paper]{amsart}
\usepackage[margin=19mm]{geometry}
\usepackage{amsmath,amssymb,mathtools}
\usepackage[scr=rsfs,cal=euler]{mathalfa}
\usepackage{xfrac}
\usepackage[unicode,hidelinks,bookmarks=false]{hyperref}
\newcommand{\ie}{i.e.\ }
\newcommand{\cf}{cf.\ }
\newcommand{\resp}{respectively\ }
\newcommand{\Subsection}{\subsection}
\providecommand{\nobreakdash}{\nobreak\hbox{-}}
\setlength{\emergencystretch}{2em}
\multlinegap0pt
\usepackage{bm}
\usepackage{bbm}
\usepackage{dsfont}
\usepackage{xspace}
\usepackage{cancel}
\usepackage{mathtools}
\usepackage{amsthm}
\makeatletter
\@ifundefined{theorem}{\newtheorem{theorem}{Theorem}}{}
\@ifundefined{lemma}{\newtheorem{lemma}[theorem]{Lemma}}{}
\makeatother
\title[Arithmetic properties of $t$-Schur overpartitions]{Arithmetic properties of $t$-Schur overpartitions}
\author{Mohammed L. Nadji, Manjil P. Saikia,, James A. Sellers}
\date{Source: arXiv:2508.17928v1 (CC BY 4.0)}
\begin{document}
\maketitle

\noindent\emph{Source and licence.} Arithmetic properties of $t$-Schur overpartitions. Version arXiv:2508.17928v1 (CC BY 4.0), distributed under the Creative Commons Attribution 4.0 International licence, as recorded in the arXiv licence metadata for this exact version and shown on the abstract page https://arxiv.org/abs/2508.17928v1. Full rights notice: licenses/arxiv-2508.17928-CC-BY-4.0.txt.\par\smallskip

\noindent\textit{Editorial scope.} Counted unit: the Theorem whose source label is \texttt{thm1} in the article (source0.tex lines 629--638, source label \texttt{thm1}), delivered together with the article's own complete original proof of it (source0.tex lines 640--718, one \texttt{proof} environment of 79 lines). No mathematics is written, repaired or re-derived: statement and proof are the article's own text line for line. Required context retained uncounted: maineq (lines 161--163); binomial (lines 208--213); binomiallemma (lines 210--212); lemma1.1 (lines 244--254); lemma1.2 (lines 244--254); lemma1.3 (lines 244--254); lemma1.4 (lines 244--254); lemma2.1 (lines 263--266); lemma4.2 (lines 280--283); lemma8.3 (lines 323--327); lemma10.1 (lines 338--340). Every definition, display and label of those lines is retained unchanged. Omitted from this package and not redistributed: 1--160, 164--207, 213--243, 255--262, 267--279, 284--322, 328--337, 341--628, 639--639, 719--1372. Nothing inside a retained excerpt is omitted or altered: every retained body is delivered exactly as the article's lines are written, so no omission receipt of any kind accompanies this package. Citations: the retained excerpts carry no citation command at all, so no bibliography entry of the article is transcribed and no reference is redistributed. Reference adaptation: none was needed and none was made; every reference inside the delivered unit resolves inside it, so no unresolved reference or citation is produced. Printed slips kept literal: no printed word, symbol, hypothesis or number of the counted statement or of its proof was altered. Layout and class: the article's own class is \texttt{amsart} with packages including hyperref, amsmath, babel, inputenc, amssymb, amsthm, graphics, graphicx, enumerate, array, bm, a4wide, color, url; the wrapper is the lane's pinned \texttt{amsart} editorial wrapper at 10pt on A4, which restates the theorem-like environments the retained text needs and adds the article's own macro definitions that the retained text needs (none), with extra wrapper packages (bm, bbm, dsfont, xspace, cancel, mathtools). The delivered numbering is the unit's own. Assets: no retained span contains a figure environment, a graphics-inclusion command, a PDF-page inclusion, a table or a TikZ picture; no image file is copied into this package, the package declares no asset dependency and no image file is redistributed with it. Licence: version arXiv:2508.17928v1 of this article is distributed under the Creative Commons Attribution 4.0 International licence, as recorded in the arXiv licence metadata for this exact version and shown on the abstract page https://arxiv.org/abs/2508.17928v1. A full rights notice is delivered at licenses/arxiv-2508.17928-CC-BY-4.0.txt. Characters: every retained excerpt is pure ASCII, so the delivered engine, XeLaTeX with its default Latin Modern text font, reports no missing character.
\begin{equation}
 \sum_{n\geq 0} \overline{S}_{t}(n)q^{n}=\prod_{n\geq 1}\dfrac{(1+q^{2n-1})(1-q^{t(2n-1)})}{(1+q^{t(2n-1)})(1-q^{2n-1})}=\dfrac{f_{2}^{3}f_{t}^{2}f_{4t}}{f_{1}^{2}f_{4}f_{2t}^{3}}=\frac{\varphi(-q^2)\varphi(-q^t)}{\varphi(-q)\varphi(-q^{2t})},  \label{maineq}
\end{equation}

\begin{lemma} \label{binomial} 
For all primes $p$ and all $k, m\geq 1$, we have 
\begin{equation}
 f_{pm}^{p^{k-1}} \equiv f_{m}^{p^{k}}  \ (\mathrm{mod} \ p^{k}). \label{binomiallemma} 
\end{equation}
\end{lemma}

\begin{align}
\frac{1}{f_1^4} &= \frac{f_4^{14}}{f_2^{14}f8^4}+4q\frac{f_4^2f_8^4}{f_2^{10}},\label{diss-2-f14}\\
\frac{f_{3}^{2}}{f_{1}^{2}} &=\frac{f_{4}^{4}f_{6}f_{12}^{2}}{f_{2}^{5}f_{8}f_{24}}+2q\frac{f_{4}f_{6}^{2}f_{8}f_{24}}{f_{2}^{4}f_{12}}, \label{lemma1.1}\\
\dfrac{f_{3}^{3}}{f_{1}} &=\dfrac{f_{4}^{3}f_{6}^{2}}{f_{2}^{2}f_{12}}+q\dfrac{f_{12}^{3}}{f_{4}}, \label{lemma1.2}\\
\frac{f_{1}}{f_{3}^{3}} &=\frac{f_{2} f_{4}^{2} f_{12}^{2}}{f_{6}^{7}} -q \frac{f_{2}^{3} f_{12}^{6}}{f_{4}^{2} f_{6}^{9}}, \label{lemma1.3}\\
\frac{f_{3}}{f_{1}^{3}} &=\frac{f_{4}^{6} f_{6}^{3}}{f_{2}^{9} f_{12}^{2}} +3q \frac{f_{4}^{2} f_{6} f_{12}^{2}}{f_{2}^{7}}, \label{lemma1.4}\\
\frac{f_{1}^{3}}{f_{3}} &=\frac{f_{4}^{3}}{f_{12}} -3q \frac{f_{2}^{2} f_{12}^{3} }{ f_{4} f_{6}^{2}}, \label{lemma1.41}\\
\dfrac{1}{f_{1}^{2} f_{3}^{2}} &= \dfrac{f_{8}^{5} f_{24}^{5}}{f_{2}^{5} f_{6}^{5} f_{16}^{2} f_{48}^{2}} +2q \dfrac{f_{4}^{4} f_{12}^{4}}{f_{2}^{6} f_{6}^{6}} +4q^{4} \dfrac{f_{4}^{2} f_{12}^{2} f_{16}^{4} f_{48}^{2}}{f_{2}^{5} f_{6}^{5} f_{8} f_{24}}, \label{lemma1.5}\\
\frac{1}{f_{1} f_{3}} &= \frac{f_{8}^{2} f_{12}^{5}}{f_{2}^{2} f_{4} f_{6}^{4} f_{24}^{2}}+q \frac{f_{4}^{5} f_{24}^{2}}{f_{2}^{4} f_{6}^{2} f_{8}^{2} f_{12}},\label{lemma1.6}\\
f_{1}f_{3} &= \frac{f_{2} f_{8}^{2} f_{12}^{4}}{f_{4}^{2} f_{6} f_{24}^{2}}-q \frac{f_{4}^{4} f_{6} f_{24}^{2}}{f_{2} f_{8}^{2} f_{12}^{2}}. \label{lemma1.7}
\end{align}

 \begin{align}
    &f_{1}f_{2}= \frac{f_{6} f_{9}^{4}}{ f_{3} f_{18}^{2}}-qf_{9}f_{18} -2q^{2} \frac{f_{3} f_{18}^{4}}{f_{6} f_{9}^{2}}, \label{lemma2.1}\\
    &\dfrac{1}{f_{1}f_{2}}=\dfrac{f_{9}^{9}}{f_{3}^{6}f_{6}^{2}f_{18}^{3}}+q\dfrac{f_{9}^{6}}{f_{3}^{5}f_{6}^{3}}+3q^{2}\dfrac{f_{9}^{3}f_{18}^{3}}{f_{3}^{4}f_{6}^{4}}-2q^{3}\dfrac{f_{18}^{6}}{f_{3}^{3}f_{6}^{5}}+4q^{4}\dfrac{f_{18}^{9}}{f_{3}^{2}f_{6}^{6}f_{9}^{3}}. \label{lemma2.2}
    \end{align}

\begin{align}
&\dfrac{f_{4}}{f_{1}}=\dfrac{f_{12}f_{18}^{4}}{f_{3}^{3}f_{36}^{2}}+ q\dfrac{f_{6}^{2}f_{9}^{3}f_{36}}{f_{3}^{4}f_{18}^{2}}+2q^{2}\dfrac{f_{6}f_{18}f_{36}}{f_{3}^{3}}, \label{lemma4.1}\\
&\frac{f_{1}}{f_{4}}= \frac{f_{6} f_{9} f_{18}}{f_{12}^{3}}-q\frac{f_{3} f_{18}^{4}}{f_{9}^{2} f_{12}^{3}}-q^{2} \frac{f_{6}^{2} f_{9} f_{36}^{3}}{f_{12}^{4} f_{18}^{2}}. \label{lemma4.2}
\end{align}

\begin{align}
\dfrac{f_{2}^{2}}{f_{1}}&=\dfrac{f_{6}f_{9}^{2}}{f_{3}f_{18}}+q\dfrac{f_{18}^{2}}{f_{9}}, \label{lemma8.1}\\
\psi(-q)&= A(q^{3}) - q \psi(-q^{9}), \label{lemma8.2}\\
&= P(-q^{3}) - q \psi(-q^{9}), \label{lemma8.3}
\end{align}

\begin{equation}
\varphi(-q)=\varphi(-q^9) -2q\dfrac{f_{3}f_{18}^{2}}{f_{6}f_{9}}. \label{lemma10.1}
\end{equation}

\begin{theorem} \label{thm1} For any nonnegative integer $n$, we have
\begin{align}
\overline{S}_{3}(9n+6) &\equiv0 \ (\mathrm{mod} \ 3), \label{thm1-cong11}\\
\overline{S}_{3}(24n) &\equiv  0 \ (\mathrm{mod} \ 3) \ \text{for} \ n>0, \label{thm1-cong1}\\
\overline{S}_{3}(24n+4) &\equiv  2 f_{1}^{4} \ (\mathrm{mod} \ 3), \label{thm1-cong2}\\
\overline{S}_{3}(24n+12) &\equiv  \psi(q) \psi(q^{3}) \ (\mathrm{mod} \ 3), \label{thm1-cong3}\\
\overline{S}_{3}(24n+16) &\equiv  0 \ (\mathrm{mod} \ 3), \label{thm1-cong4}\\
\overline{S}_{3}(27n) &\equiv  \overline{S}_{3}(3n) \ (\mathrm{mod} \ 3). \label{thm1-cong5}
\end{align}
\end{theorem}

\begin{proof} Setting $t=3$ in \eqref{maineq} and then substituting \eqref{lemma1.1} into the resulting equation, we find that
\[
\sum_{n \geq 0} \overline{S}_{3}(n) q^{n} = \dfrac{f_{4}^{3} f_{12}^{3}}{f_{2}^{2} f_{6}^{2} f_{8} f_{24}} + 2q \dfrac{f_{8} f_{24}}{f_{2} f_{6}}.
\]
Equating the even and odd powers on both sides of the above equation, we get
\begin{align}
&\sum_{n \geq 0} \overline{S}_{3}(2n) q^{n} = \dfrac{f_{2}^{3} f_{6}^{3}}{f_{1}^{2} f_{3}^{2} f_{4} f_{12}}, \label{a1}\\
&\sum_{n \geq 0} \overline{S}_{3}(2n+1) q^{n} = 2\dfrac{f_{4} f_{12}}{f_{1} f_{3}}. \label{a2}
\end{align}
Applying Lemma \ref{binomial} to \eqref{a1}, with $p=3$ and $k=1$, we obtain
\begin{equation}
\sum_{n \geq 0} \overline{S}_{3}(2n) q^{n} \equiv  \dfrac{f_{1} f_{2}^{3} f_{6}^{3}} {f_{3}^{3} f_{4} f_{12}} \ (\mathrm{mod} \ 3). \label{a3}
\end{equation}
Using \eqref{lemma1.3} in \eqref{a3} and then extracting the terms of the form $q^{2n}$ from both sides of the resulting equation, we arrive at
\begin{equation}
\sum_{n \geq 0} \overline{S}_{3}(4n) q^{n} \equiv  \dfrac{f_{1}^{4} f_{2} f_{6}} {f_{3}^{4}} \equiv \dfrac{f_{1} f_{2} f_{6}} {f_{3}^{3}}  \ (\mathrm{mod} \ 3). \label{a4}
\end{equation}
Substituting \eqref{lemma2.1} into \eqref{a4}, we find that
\begin{equation}
\sum_{n \geq 0} \overline{S}_{3}(4n) q^{n} \equiv  \dfrac{f_{6}^{2} f_{9}^{4}} {f_{3}^{4} f_{18}^{2}} - q \dfrac{f_{6} f_ {9} f_{18}}{f_{3}^{3}} -2q^{2} \dfrac{f_{18}^{4}}{f_{3}^{2} f_{9}^{2}}  \ (\mathrm{mod} \ 3). \label{a5}
\end{equation}
Extracting the terms of the form $q^{3n+j}$ for $j=0,1$  from both sides of equation \eqref{a5}, we arrive at
\begin{align}
&\sum_{n \geq 0} \overline{S}_{3}(12n) q^{n} \equiv  \dfrac{f_{2}^{2} f_{3}^{4}} {f_{1}^{4} f_{6}^{2}} \equiv \dfrac{f_{2}^{2} f_{3}^{3}} {f_{1} f_{6}^{2}}   \ (\mathrm{mod} \ 3), \label{a6}\\
&\sum_{n \geq 0} \overline{S}_{3}(12n+4) q^{n} \equiv  2 \dfrac{f_{2} f_ {3} f_{6}}{f_{1}^{3}}  \ (\mathrm{mod} \ 3). \label{a7}
\end{align}
Employing \eqref{lemma1.2} into \eqref{a6} and then extracting the even and the odd powers from both sides of the resulting equation, we find that
\begin{align}
&\sum_{n \geq 0} \overline{S}_{3}(24n) q^{n} \equiv  \dfrac{f_{2}^{3}} {f_{6}} \equiv 1 \ (\mathrm{mod} \ 3), \label{a8}\\
&\sum_{n \geq 0} \overline{S}_{3}(24n+12) q^{n} \equiv   \dfrac{f_{1}^{2} f_ {6}^{3}}{f_{2} f_  {3}^{2}} \equiv \dfrac{f_{2}^{2} f_ {6}^{2}}{f_{1} f_{3}}  \ (\mathrm{mod} \ 3). \label{a9}
\end{align}
Congruence \eqref{thm1-cong1} follows from \eqref{a8} and congruence \eqref{thm1-cong3} follows from \eqref{a9} by using the definition of $\psi(q)$.

Substituting \eqref{lemma1.4} into \eqref{a7}, we get
\begin{equation}
\sum_{n \geq 0} \overline{S}_{3}(12n+4) q^{n} \equiv  2 \dfrac{f_{4}^{6} f_{6}^{4}}{f_{2}^{8} f_{12}^{2}} + 6 q^{2}  \dfrac{f_{4}^{2} f_{6}^{2} f_{12}^{2}}{f_{2}^{6}}  \ (\mathrm{mod} \ 3).  \label{a10}
\end{equation}
Congruence \eqref{thm1-cong4} follows from the above equation by extracting the odd powers from both sides. Equating the even powers from both sides of equation \eqref{a10}, we find that
\begin{equation}
\sum_{n \geq 0} \overline{S}_{3}(24n+4) q^{n} \equiv  2 \dfrac{f_{2}^{6} f_{3}^{4}}{f_{1}^{8} f_{6}^{2}} \equiv 2f_{1}^{4}  \ (\mathrm{mod} \ 3).  \label{a11}
\end{equation}
Congruence \eqref{thm1-cong2} is an immediate result from \eqref{a11}.

Again, setting $t=3$ in \eqref{maineq} and applying \eqref{binomiallemma}, with $p=3$ and $k=1$, we get 
\begin{equation*}
\sum_{n \geq 0} \overline{S}_{3}(n) q^{n} = \frac{f_{2}^3 }{f_{1}^2 f_{4}} \cdot \frac{f_{3}^2 f_{12}}{f_{6}^3} \equiv \frac{f_{1} f_{3} f_{12}}{f_{4} f_{6}^2} \equiv \frac{f_{1}}{f_{4}} \cdot  \frac{f_{3} f_{12}}{f_{6}^2} \ (\mathrm{mod} \ 3).
\end{equation*}
Now by invoking \eqref{lemma4.2} into the above equation and then extracting the terms of the form $q^{3n}$ from both sides of the resulting equation, we arrive at
\begin{equation*}
\sum_{n \geq 0} \overline{S}_{3}(3n) q^{n} \equiv \frac{f_1 f_3 f_6}{f_2 f_4^2} \equiv \frac{f_1 f_3 f_4 f_6}{f_2 f_{12}} = \psi(-q) \frac{f_3 f_6}{f_{12}} \ (\mathrm{mod} \ 3).
\end{equation*}
By using \eqref{lemma8.3} in the above equation, we see that
\begin{equation}
\sum_{n \geq 0} \overline{S}_{3}(3n) q^{n} \equiv\big( P(-q^3) - q \psi(-q^9) \big) \frac{f_3 f_6}{f_{12}} \ (\mathrm{mod} \ 3),  \label{sel01}
\end{equation}
and since the above $3$-dissection contains no terms of the form $q^{3n+2}$, we deduce that
\begin{equation*}
\overline{S}_{3}(9n+6)  \equiv 0 \ (\mathrm{mod} \ 3).
\end{equation*}

From \eqref{sel01}, we know that
\begin{align*}
\sum_{n \geq 0} \overline{S}_{3}(9n) q^{n} &\equiv P(-q) \frac{f_1 f_2}{f_{4}} = \frac{f_1 f_4 f_6^5}{f_2^2 f_3^2 f_{12}^2}  \cdot \frac{f_1 f_2}{f_{4}} = \frac{f_1^2 f_6^5}{f_2 f_3^2 f_{12}^2} \\
&\equiv  \varphi(-q) \bigg( \frac{ f_6^5}{ f_3^2 f_{12}^2} \bigg)
\ (\mathrm{mod} \ 3).
\end{align*}
By using \eqref{lemma10.1} in the above equation, we arrive at
\[
\sum_{n \geq 0} \overline{S}_{3}(9n) q^{n} \equiv \bigg(  \varphi(-q^9) -2q\dfrac{f_{3}f_{18}^{2}}{f_{6}f_{9}} \bigg) \bigg( \frac{ f_6^5}{ f_3^2 f_{12}^2} \bigg) \ (\mathrm{mod} \ 3).
\]
Thus, 
\[
\sum_{n \geq 0} \overline{S}_{3}(27n) q^{n} \equiv  \varphi(-q^3) \frac{ f_2^5}{ f_1^2 f_{4}^2} =  \frac{ f_2^5 f_3^2}{ f_1^2 f_{4}^2 f_6} \equiv \frac{f_1 f_3 f_6}{f_2 f_4^2} \equiv \sum_{n \geq 0} \overline{S}_{3}(3n) q^{n}  \ (\mathrm{mod} \ 3).
\]
Therefore, for all $n\geq 0$, we have
\[
\overline{S}_{3}(27n) \equiv  \overline{S}_{3}(3n) \ (\mathrm{mod} \ 3).
\]
\end{proof}

\end{document}
